
% ==================================================
% 2. PROJECT OVERVIEW
% ==================================================

\section{Project Overview}


% ==================================================
% 2.1 PROJECT PURPOSE, SCOPE, AND OBJECTIVES
% ==================================================

\subsection{Project Purpose, Scope, and Objectives}

Many patients, especially those with chronic diseases, take
multiple medications prescribed by different doctors.
These medications may interact with each other or contain
the same active ingredients, leading to potential health
risks. Different brand names and unclear prescriptions
make it difficult for patients to identify these risks
and manage their medication schedules.

The Medication Safety Assistant aims to help patients
identify potential medication risks and manage their
prescribed medication schedules.
Users can upload prescription images, type natural-language
requests, use voice input, or enter medication details
through manual forms.



\PAChange{The system checks supported drug--drug
interactions and duplicate active ingredients
when sufficient data is available. It provides
safety warnings, clear explanations, and
recommended next steps, and supports simulated
healthcare professional (HCP) review when needed.}


It also offers medication reminders and adherence tracking
to help patients follow their confirmed schedules.

\PAChange{Supporting features include medication
and safety history, an AI chatbot, patient health
summaries, medication dataset review, user and
role management, and technical support.}

The project objectives are:

\begin{itemize}

    \item Support medication input through OCR,
    natural-language text, voice, and manual forms,
    with user confirmation before assessment.

    \item Identify user intent and extract relevant
    medication information from natural-language text.

    \item \PAChange{Identify Vietnamese medications
    and active ingredients using permitted reference
    data, and perform supported drug--drug interaction
    checks using usable DDInter 2.0 data and verified
    ingredient mappings.}

    \item \PAChange{Detect duplicate active ingredients
    when sufficient data is available. Report detected
    risks, severity, assessment status, and data
    limitations without treating unknown results
    as safe.}

    \item \PAChange{Support simulated HCP review
    of medication safety cases, including patient
    permission, follow-up questions, case-based
    conversations, and recorded review decisions.}

    \item \PAChange{Provide medication reminders
    based on verified prescription schedules or
    recorded HCP instructions. Apply the severe
    interaction safety gate and track dose status
    and adherence.}

    \item \PAChange{Allow patients to view medication
    and safety history, ask supported medication
    questions, and generate and share reviewed
    health summaries with authorized HCPs.}

    \item \PAChange{Support medication dataset review,
    user authentication, role-based access control,
    administrative activity history, and technical
    issue management.}

    \item \PAChange{Test the main workflows, safety
    restrictions, missing-data handling, and
    access permissions using controlled test cases.}

    \item Deliver a working Web prototype within
    five sprints.

\end{itemize}

\PAChange{The prototype is developed for research
and educational use. HCP review is simulated
and does not represent real clinical approval.
The system does not independently diagnose,
prescribe, stop, or change medication treatment.}

\PAChange{The current scope excludes ADR reporting,
drug--food interaction checks, advanced medication
timing analysis, pill recognition, handwritten
prescription support, Personal Medication Inventory,
and integration with real healthcare systems.}

\PAChange{Consultation preparation, medication
record comparison, general disease information,
and emergency guidance are also outside the
current prototype scope. These features may
be considered in future research.}


% ==================================================
% 2.2 ASSUMPTIONS AND CONSTRAINTS
% ==================================================

\subsection{Assumptions and Constraints}

\begin{longtable}{|
    L{0.18\textwidth}|
    L{0.35\textwidth}|
    L{0.35\textwidth}|
}
\hline

\textbf{Type} &
\textbf{Description} &
\textbf{Effect on the project}
\\ \hline

\endfirsthead

\hline
\textbf{Type} &
\textbf{Description} &
\textbf{Effect on the project}
\\ \hline

\endhead


% --------------------------------------------------
% USER CONFIRMATION
% --------------------------------------------------

Assumption
& \PAChange{Users can review, correct, and confirm
recorded medication information from OCR,
NLP, Voice, or Manual Input.}
& \PAChange{Unconfirmed information cannot be
used for safety assessment or reminder
activation. User confirmation does not
prove medication identity or medical safety.}
\\ \hline


% --------------------------------------------------
% MEDICATION DATA
% --------------------------------------------------

Assumption
& \PAChange{Permitted Vietnamese medication
reference data is available for
research and prototype testing.}
& \PAChange{The team must verify source rights,
data coverage, and ingredient mappings.
Unsupported medication identities
remain unresolved.}
\\ \hline


% --------------------------------------------------
% DRUG INTERACTION DATA
% --------------------------------------------------

\PAChange{Data dependency}
& \PAChange{DDInter 2.0 safety checks depend on
usable interaction data and verified
ingredient mappings.}
& \PAChange{Unsupported medicine pairs remain
Unassessed. Other supported checks may
continue without claiming full safety.}
\\ \hline


% --------------------------------------------------
% DATA READINESS
% --------------------------------------------------

\PAChange{Data readiness constraint}
& \PAChange{The medication dataset has undergone
automated validation, but accepted product
records do not guarantee verified
drug-interaction mappings.}
& \PAChange{DDI coverage must be checked
separately. The team will use controlled
test cases when required mappings are
unavailable and report assessment limits.}
\\ \hline


% --------------------------------------------------
% SIMULATED HCP REVIEW
% --------------------------------------------------

Assumption
& \PAChange{Authorized HCP accounts are available
for simulated case review and patient--HCP
conversations.}
& \PAChange{Review requests remain Pending when
no reviewer is available. Simulated decisions
do not represent actual clinical approval.}
\\ \hline


% --------------------------------------------------
% SCHEDULE
% --------------------------------------------------

Schedule constraint
& The project has three members and five
two-week sprints.
& \PAChange{Tasks must be prioritized according
to dependencies, safety risks, and available
team capacity to meet the deadline.}
\\ \hline


% --------------------------------------------------
% RESOURCES
% --------------------------------------------------

Resource constraint
& Development uses student laptops, available datasets,
and free or approved tools.
& \PAChange{The team must account for limited
computing resources, development time,
and external service availability.}
\\ \hline


% --------------------------------------------------
% INPUT PROCESSING
% --------------------------------------------------

Technical constraint
& OCR and Voice may produce incorrect results.
Handwritten prescriptions are not supported.
& Users must be able to correct errors or enter medication
information manually.
\\ \hline


% --------------------------------------------------
% INCOMPLETE SAFETY DATA
% --------------------------------------------------

Data constraint
& \PAChange{Medication identities, active-ingredient
mappings, or interaction data may be incomplete
or unavailable.}
& \PAChange{The system must distinguish Assessed,
Partially Assessed, Unassessed, and Failed results.
Missing information must never be presented
as proof of safety.}
\\ \hline


% --------------------------------------------------
% MEDICATION SAFETY
% --------------------------------------------------

Safety constraint
& The system cannot independently prescribe, stop,
or change medications.
& \PAChange{Warnings must use supported findings.
Detected severe interactions must be reported
immediately without automatically stopping
prescribed treatment or existing reminders.}
\\ \hline


% --------------------------------------------------
% MEDICATION REMINDERS
% --------------------------------------------------

Reminder constraint
& \PAChange{Reminder schedules must come from
existing prescriptions or recorded HCP
instructions and be verified by users.
Proposed changes to prescribed timing
require simulated HCP confirmation.}
& \PAChange{Detected severe interactions hold
new reminders for affected medicines
pending required review. Unrelated eligible
reminders may proceed. Existing reminders
are not automatically stopped.}
\\ \hline


% --------------------------------------------------
% PRIVACY
% --------------------------------------------------

Privacy constraint
& Medication records, prescription images, and voice
data may contain sensitive information.
& \PAChange{The prototype uses synthetic or
de-identified test data, limits collection,
and restricts access to authorized users
and information.}
\\ \hline


% --------------------------------------------------
% USER ACCESS
% --------------------------------------------------

\PAChange{Access constraint}
& \PAChange{Five user roles have different
permissions. Patient information may
be shared only with authorized HCPs
after patient permission.}
& \PAChange{Role-Based Access Control
and case-level permissions must
be implemented and tested.
Unauthorized access must be denied.}
\\ \hline


% --------------------------------------------------
% PLATFORM
% --------------------------------------------------

\PAChange{Platform dependency}
& \PAChange{Camera, microphone, and browser
notifications depend on device support
and user permissions.}
& \PAChange{Manual medication input remains
available. Reminder schedules remain
viewable when notifications are unavailable.}
\\ \hline


% --------------------------------------------------
% EXTERNAL DATA RIGHTS
% --------------------------------------------------

\PAChange{External constraint}
& \PAChange{Medication datasets and external
resources may have usage or redistribution
restrictions.}
& \PAChange{The team must review applicable
source terms before storing, sharing,
or publishing datasets in the repository.}
\\ \hline

\end{longtable}


% ==================================================
% 2.3 PROJECT DELIVERABLES
% ==================================================

\subsection{Project Deliverables}

\PAChange{The following deliverables define the
expected outputs of each development stage.
Acceptance conditions will be checked through
document reviews, demonstrations, or
controlled tests as appropriate.}

\begin{longtable}{|
    L{0.24\textwidth}|
    L{0.41\textwidth}|
    L{0.13\textwidth}|
    L{0.12\textwidth}|
}
\hline

\textbf{Deliverable} &
\textbf{Description / Acceptance Condition} &
\textbf{Owner} &
\textbf{Target}
\\ \hline

\endfirsthead

\hline
\textbf{Deliverable} &
\textbf{Description / Acceptance Condition} &
\textbf{Owner} &
\textbf{Target}
\\ \hline

\endhead


% --------------------------------------------------
% PA1
% --------------------------------------------------

PA1 planning package
& Initial Project Plan, Vision Document,
3--6 UI screens, two weekly reports,
and accessible tracking evidence.
& Whole team
& Sprint 1
\\ \hline


% --------------------------------------------------
% PA2 REQUIREMENTS PACKAGE
% --------------------------------------------------

\PAChange{PA2 requirements package}
& \PAChange{Revised Vision Document (13 PF,
10 NFR), revised Project Plan,
18 Use-case Specifications, corresponding
Use-case Diagrams, and an updated UI prototype.
Documents must be consistent and reviewed
against PA2 requirements.}
& \PAChange{Whole team}
& \PAChange{Sprint 2}
\\ \hline


% --------------------------------------------------
% ARCHITECTURE AND UX
% --------------------------------------------------

\PAChange{Architecture and UX baseline}
& \PAChange{System context, main workflows,
module interfaces, common medication data
format, and UI designs for key user journeys.
Major integration decisions are documented
and reviewed by the team.}
& \PAChange{Trang /
Team}
& \PAChange{Sprint 2}
\\ \hline


% --------------------------------------------------
% MEDICATION DATA BASELINE
% --------------------------------------------------

\PAChange{Medication data baseline}
& \PAChange{Prepared medication records,
ingredient information, dataset audit results,
source references, and documented coverage.
DDI mapping limitations and data-use
restrictions are recorded.}
& \PAChange{Long}
& \PAChange{Sprint 2--3}
\\ \hline


% --------------------------------------------------
% MEDICATION INTAKE POC
% --------------------------------------------------

Medication Intake Proof of Concept
& \PAChange{OCR, NLP, Voice/STT, and Manual
Input produce editable medication information
in a common format. Users can confirm
recorded details or recover from
input-processing failures.}
& Whole team
& Sprint 3
\\ \hline


% --------------------------------------------------
% SAFETY ASSESSMENT POC
% --------------------------------------------------

Medication Safety Proof of Concept
& Prepared Vietnamese medication data, reviewed ingredient
mappings to DDInter 2.0, and tested drug--drug interaction
and duplicate-ingredient checks with clear risk and
incomplete-data results.
& Long & Sprint 3\\ \hline

Integrated test build
& \PAChange{Integrate medication intake,
identification, safety checks, result display,
and simulated HCP review. Test case sharing,
patient--HCP messages, review decisions,
and major error paths.}
& Whole team
& Sprint 4
\\ \hline


% --------------------------------------------------
% PATIENT SUPPORT FEATURES
% --------------------------------------------------

\PAChange{Patient support workflows}
& \PAChange{Provide medication and safety
history, supported chatbot questions,
patient health summaries, and
permission-based summary sharing.
Test access and information accuracy
within available data limitations.}
& \PAChange{Long /
Team}
& \PAChange{Sprint 4}
\\ \hline


% --------------------------------------------------
% SUPPORTING SYSTEM WORKFLOWS
% --------------------------------------------------

\PAChange{Supporting system workflows}
& \PAChange{Provide dataset review,
authentication, role and account management,
administrative activity history,
and technical issue reporting and handling.
Test role restrictions and recorded actions.}
& \PAChange{Whole team}
& \PAChange{Sprint 4--5}
\\ \hline


% --------------------------------------------------
% FINAL PROTOTYPE
% --------------------------------------------------

\PAChange{Final prototype release}
& \PAChange{Integrated Web research prototype
covering the agreed current-scope workflows,
verified-schedule reminders, severe DDI
safety gate, dose status, and adherence history.
Complete regression tests, resolve critical
defects, and prepare documentation and demo
materials.}
& Whole team
& Sprint 5
\\ \hline


% --------------------------------------------------
% PROJECT EVIDENCE
% --------------------------------------------------

Project evidence repository
& Versioned source, decisions, test evidence,
sprint records, risk register, meeting notes,
and weekly reports accessible to teaching staff.
& Trang
& Weekly
\\ \hline

\end{longtable}
